%! suppress = MissingLabel
%! suppress = TooLargeSection
%! suppress = EscapeUnderscore
%%
%% Build with xelatex + biber (the engine the official ISP RAS template and its
%% Makefile use; the Russian bibliography heading needs Unicode fonts, so the
%% pdflatex/T1 path fails on it):
%%   latexmk -pdf -xelatex effect-systems-free-variables.tex
%%
\documentclass{ProcISPRAS}

% Page geometry as in the official ISP RAS template.
\usepackage[papersize={14.86cm,21cm},
            left=1.5cm, right=1cm, top=0.8cm, bottom=1cm,
            includehead, headheight=8pt, heightrounded, headsep=6pt,
            includefoot, footskip=16pt]{geometry}
\usepackage{minted}

\addbibresource{bib.bib}

\usepackage{url}
\usepackage{amsmath,amssymb,amsfonts}
\usepackage{proof}
\usepackage{booktabs}
\usepackage{array}
\usepackage{listings}
\usepackage{tikz}
\usetikzlibrary{positioning, arrows.meta}

% --- math shorthands reused from the original overview ---
\newcommand{\seq}{;~}
\newcommand{\ap}{~}

% --- listings: lightweight definitions for the languages used inline ---
\lstdefinestyle{modern}{
    basicstyle=\ttfamily\footnotesize,
    keywordstyle=\bfseries,
    commentstyle=\itshape,
    stringstyle=,
    numberstyle=\tiny,
    tabsize=2,
    showspaces=false,
    showstringspaces=false,
    breaklines=true,
    breakatwhitespace=true,
    captionpos=b,
    escapeinside={(*@}{@*)},
    columns=fullflexible,
    keepspaces=true,
}
\lstset{style=modern}
\lstdefinelanguage{koka}{
    keywords={fun,val,var,if,else,match,with,return,do,handle,handler,ctl,throw,try,effect,rec,mutable,fn,mask,resume},
    sensitive=true, comment=[l]{//}, morecomment=[s]{/*}{*/},
    morestring=[b]", morestring=[b]', alsoletter={!},
}
\lstdefinelanguage{effekt}{
    keywords={def,val,var,if,else,try,with,box,unbox,interface,effect,handle,resume,at},
    sensitive=true, comment=[l]{//}, morecomment=[s]{/*}{*/}, morestring=[b]",
}
\lstdefinelanguage{colang}{
    keywords={let,var,effect,perform,data,context,fun,where,if,else,throw,try,catch,local,free,op,handle,resume,return},
    sensitive=true, comment=[l]{//}, morecomment=[s]{/*}{*/}, morestring=[b]",
}

% =====================================================================
%% Header running titles
\titleheadru{Системы эффектов как управление свободными переменными}
\titleheaden{Effect systems as management of free variables}

\volhead{37}
\issuehead{1}
\pageshead{1--24}
\yearhead{2026}
\setcounter{page}{1}

% --- authors (bilingual) ---
\newauthor
\inst{1}
\authorname{A.\,S. Stoyan}{А.\,С. Стоян}
\orcid{0009-0004-6935-1447}
\email{a.stoyan@hse.ru}

\affil[1]{
National Research University Higher School of Economics,\\
3A, Kantemirovskaya Str., Saint Petersburg, 194100, Russia}{
Национальный исследовательский университет <<Высшая школа экономики>>,\\
194100, Россия, г.~Санкт-Петербург, ул.~Кантемировская, д.~3А}

\title{Effect Systems as a Mechanism for Managing Free Variables: A Unifying Perspective}{%
Системы эффектов как механизм управления свободными переменными: единый взгляд}

\doi{10.15514/ISPRAS-2026-XX(X)-XX} % TODO: replace with the DOI assigned on acceptance

\keywords{%
effect systems; computational effects; free variables; implicit parameters;
escape analysis; capabilities; effect rows; modal types; type systems;
contextual polymorphism}{%
системы эффектов; вычислительные эффекты; свободные переменные; неявные параметры;
escape-анализ; capabilities; строки эффектов; модальные типы; системы типов;
контекстный полиморфизм}

\abstract{% English abstract
Effect systems lift a program's computational effects to the type level, sharpening
abstraction boundaries and letting the compiler check that every computation runs in a
supporting context. Yet they are rarely adopted in mainstream languages. One reason, we
believe, is that the leading designs --- row-polymorphic types, capabilities, and modal
types --- are usually presented as unrelated and alien approaches.

We offer a unifying
reading: an effect system manages the \emph{free} variables of a computation, whereas a
classical type system governs its \emph{bound} variables. A contextual requirement passes
through a lifecycle --- it arises \emph{unfulfilled}, is \emph{fulfilled} by a value from
the context (an implicit argument), is \emph{captured} by a closure and held in place by
type-based escape analysis, and may be released back into an unfulfilled state at a context
boundary --- and the families of effect systems differ only in which phase they make
explicit in types. Row-based systems (Koka, Links, Java checked exceptions, the Haskell mtl
idiom) make the unfulfilled phase explicit; capability-based systems (Effekt, Scala capture
checking, second-class values) make the captured phase explicit and prevent escape; modal
systems mark the transitions between phases, controlling leakage rather than forbidding it.
The perspective suggests a constructive conclusion: an effect system can be assembled from
two ingredients already familiar to practitioners --- implicit parameters and type-based
escape analysis.
}{% Russian abstract
Системы эффектов поднимают сведения о вычислительных эффектах программы на уровень типов,
уточняя границы абстракций и позволяя компилятору проверять, что вычисление выполняется в
поддерживающем его контексте. Тем не менее в массовых языках они применяются редко. Одна из
причин, на наш взгляд, в том, что ведущие подходы --- полиморфизм по строкам, capabilities и
модальные типы --- обычно подаются как несвязанные и чуждые привычной практике.

Мы предлагаем
объединяющий взгляд: система эффектов управляет \emph{свободными} переменными вычисления,
тогда как классическая система типов --- его \emph{связанными} переменными. Контекстное
требование проходит жизненный цикл: оно возникает \emph{невыполненным}, \emph{выполняется}
значением из контекста (неявным аргументом), \emph{захватывается} замыканием и удерживается
на месте escape-анализом на уровне типов, а на границе контекста может быть возвращено в
невыполненное состояние; семейства систем эффектов различаются лишь тем, какую фазу этого
цикла они делают явной в типах. Строковые системы (Koka, Links, проверяемые исключения Java,
идиома mtl в Haskell) делают явной фазу невыполненного требования; системы на основе
capabilities (Effekt, отслеживание захвата в Scala, значения второго класса) делают явной
фазу захвата и предотвращают утечку; модальные системы помечают переходы между фазами,
управляя утечкой, а не запрещая её. Этот взгляд подсказывает конструктивный вывод: систему
эффектов можно собрать из двух знакомых практикам возможностей --- неявных параметров и
escape-анализа на уровне типов.
}

\begin{document}
\selectlanguage{russian}

\makedoi
\maketitleru

\section{Введение в вычислительные эффекты}

Программирование в значительной степени сводится к управлению сложностью:
разработчик стремится сосредоточиться на отдельном фрагменте поведения, не держа
в голове всю программу целиком. Абстракция и разделение ответственности --- ключевые
приёмы: различные обязанности делегируются разным сущностям, таким как функции или
модули. Часто в качестве такой сущности удобно выделять \emph{контекст исполнения}
вычисления, а \emph{эффекты} возникают естественным образом как взаимодействия
вычисления с этим контекстом~\cite{kiselyov2013extensible}.

Начнём с противоположного понятия --- чистоты. Единственный наблюдаемый результат
чистого вычисления есть его значение. Функцию называют чистой, если она возвращает
одинаковый результат на одинаковых аргументах, а её применение является чистым
вычислением. Программирование с чистыми функциями обычно считается хорошей
практикой: композиция чистых функций чиста, семантика кода прозрачна, а системы
типов работают хорошо, обеспечивая частичную спецификацию, безопасность и ясность
абстракций.

Однако писать и сопровождать нетривиальную программу, пользуясь лишь чистыми
вычислениями, трудно. Даже ввод-вывод можно смоделировать явной передачей значения
\lstinline|World| по программе~\cite{peyton1993imperative}: тогда всё приходится
передавать в параметрах и возвращать в результатах, а программист вынужден вручную
вести учёт по всему коду. Нужен инструмент управления сложностью --- нечто, что
прячет несущественные детали.

\begin{lstlisting}[language=haskell]
    getList :: Int -> World -> (World, [Int])
    getList n w | n == 0 = (w, [])
                | otherwise =
      let (w', x) = getInt w in
      let (w'', xs) = getList (n - 1) w' in
      (w'', x : xs)
\end{lstlisting}

Вычислительные эффекты и есть такой инструмент. Эффект --- это взаимодействие
вычисления с контекстом выполнения, ответственным за учёт; вычисление обращается к
контексту напрямую, без церемоний, которыми полон чистый код. Контекстом может
служить система времени выполнения, предоставляющая сервисы ОС, изменяемые ячейки
памяти, развитые механизмы управления, ввод-вывод и т.\,п. \emph{Эффектная операция}
--- это конструкция, порождающая эффект. Например, монадическая машинерия Haskell и
\lstinline|do|-нотация позволяют абстрагироваться от передачи мира.

\begin{lstlisting}[language=haskell]
    getList :: Int -> IO [Int]
    getList n | n == 0 = pure []
              | otherwise = do
      x <- getInt
      xs <- getList (n - 1)
      pure (x : xs)
\end{lstlisting}

Некоторые подходы позволяют программисту определять собственные эффекты и контексты
выполнения как библиотечный код. В частности, \emph{обработчики эффектов} (effect
handlers) --- многообещающий подход, зарекомендовавший себя и за пределами
сообщества чисто функциональных языков~\cite{plotkin2013handling,
chandrasekaran2018algebraic}. Обработчики разделяют сигнатуру эффекта и его
реализацию~\cite{kiselyov2013extensible}: выражение можно рассматривать как клиента,
который вычисляется либо в чистый результат, либо в запрос к серверу (контексту
выполнения, обработчику эффекта). Обработчик задаёт область, в которой возможны
определённые эффектные операции, и интерпретирует их как чистое значение и эффекты
объемлющей области. Собственные эффекты служат инструментом построения встроенных
предметно-ориентированных языков~\cite{leijen2017structured}, а различные модели
эффектов соревнуются в выразительности, производительности и удобстве композиции
независимо определённых эффектов~\cite{liang1995monad, kiselyov2013extensible,
schrijvers2019monad, van2024framework}.

Многие эффектные операции допустимы лишь в ограниченной области. Например,
исключение можно бросить только в контексте соответствующего обработчика, а
изменяемую ячейку, размещённую на стеке, нельзя использовать после завершения
функции. Статический контроль над применением таких операций предотвращает множество
ошибок времени выполнения и повышает надёжность программ. \emph{Системы эффектов} как раз
и являются механизмом обеспечения такого контроля на уровне типов. Также, так как эффекты по своей природе неявны на уровне термов, об их наличии сложно рассуждать по исходному коду.
Поэтому всё же требуется некоторая степень явности, которую системы эффектов обеспечивают
на уровне типов. Это помогает задавать чёткие функциональные абстракции и не давать деталям
реализации незаметно утекать в виде взаимодействий с контекстом.

Было предложено множество перспективных конструкций систем эффектов:
полиморфные по строкам типы~\cite{leijen2014koka}, capabilities~\cite{
brachthauser2022effects, boruch2023capturing}, модальные типы~\cite{tang2025modal}
и др. Однако каждый подход обладает своими достоинствами и недостатками, а общая
картина пространства проектных решений изучена недостаточно. При этом проектировщику
языка важно иметь целостное представление о доступных вариантах, чтобы выбрать
наиболее подходящее решение с учётом особенностей и ограничений целевого языка.

В настоящей работе мы предлагаем единый взгляд на эти, казалось бы, разрозненные
подходы: \emph{система эффектов --- это механизм управления свободными переменными
вычисления на уровне типов}. Опираясь на наблюдения из языка
Scala~\cite{odersky2021safer, boruch2023capturing}, этот взгляд позволяет
рассматривать классическую систему типов и систему эффектов как две стороны одной
задачи: первая управляет \emph{связанными} переменными, вторая --- \emph{свободными}.
Свободные переменные, в свою очередь, проходят единый \emph{жизненный цикл контекстного
требования}, и каждое из существующих семейств систем эффектов оказывается выбором
того, какую фазу этого цикла делать явной в типах. Настоящая работа --- именно
объединяющая перспектива, а не исчерпывающий обзор: охват сознательно сужен до трёх
семейств. Она основана на нашем более раннем англоязычном обзоре~\cite{stoyan2024modern},
но отличается от него по существу: материал реорганизован вокруг тезиса о свободных
переменных, каждое семейство снабжено отдельным прочтением с этих позиций
(разделы~\ref{subsec:rows-reading}, \ref{subsec:caps-reading}
и~\ref{subsec:modal-reading}), а сравнение проведено вдоль единой оси фаз
(раздел~\ref{sec:comparison}). Из этого взгляда вытекает и конструктивная программа:
систему эффектов можно собрать из неявных параметров и escape-анализа на уровне типов ---
эта идея служит качественной рамкой всего изложения.

Вклад работы можно резюмировать так:
\begin{itemize}
    \item мы напоминаем общие свойства и желаемые характеристики систем эффектов,
    фиксируя терминологию для дальнейшего изложения (раздел~\ref{sec:general});
    \item мы формулируем единую точку зрения на системы эффектов как на средство
    управления свободными переменными и описываем \emph{жизненный цикл контекстного
    требования}, вдоль фаз которого располагаются все рассматриваемые конструкции
    (раздел~\ref{sec:idea});
    \item мы заново рассматриваем три семейства современных систем эффектов ---
    строковые (раздел~\ref{sec:rows}), на основе capabilities
    (раздел~\ref{sec:capabilities}) и модальные (раздел~\ref{sec:modal}) ---
    и показываем, какую фазу этого цикла делает явной каждое из них;
    \item мы проводим сравнительный анализ семейств вдоль единой оси фаз и показываем,
    что систему эффектов можно собрать как сочетание двух знакомых практикам
    возможностей --- неявных параметров и escape-анализа на уровне типов
    (раздел~\ref{sec:comparison}).
\end{itemize}


\section{Системы эффектов в целом}
\label{sec:general}

В этом разделе описывается базовая идея, лежащая в основе всех систем эффектов, и
очерчиваются их желаемые свойства, к которым мы будем возвращаться на протяжении всей
работы. Для пояснения общих идей мы
пользуемся неформальной нотацией на основе System~F~\cite{girard1971extension},
расширяя её по мере необходимости. Рассматриваются только языки с вызовом по значению,
поскольку результаты обычно адаптируются к монадической
постановке~\cite{wadler2003marriage}; тогда эффекты возникают лишь при применении
функций.

Классические системы типов контролируют форму входных аргументов и результата. Так,
функция отображает натуральные в натуральные: $\lambda x\ldotp x + 1 : nat\to nat$.
Однако, как обсуждалось, важно отслеживать не только чистый аспект вычислений, но и
то, как функция взаимодействует с контекстом выполнения при работе, то есть какие
эффекты она может произвести. Метку эффекта можно прикрепить к самой стрелке:
\[\lambda x\ldotp print \ap x\seq x + 1 : nat\to^{print} nat\]
Система эффектов отслеживает эффекты для каждого терма, но в типе фиксирует лишь
эффекты абстракций. Дело в том, что в языке с вызовом по значению абстракция ---
единственный способ отложить вычисление; такое отложенное вычисление можно передавать,
и оно потенциально покидает исходный контекст. Поэтому анализ становится
чувствительным к потоку данных, и эффекты удобнее отслеживать явно. Как мы покажем в
разделе~\ref{sec:idea}, явного отслеживания
требуют именно те свободные переменные, которые могут утечь вместе с отложенным
вычислением.

Очертим желаемые свойства систем эффектов; многие системы жертвуют частью из них ради
простоты.

\textbf{Безопасность эффектов.} Всякий эффект может выполняться лишь в контексте,
который его поддерживает; иначе систему эффектов следует считать обязанной отвергнуть
программу. В частности, всякое исключение должно быть поймано, а в общем случае всякий
эффект должен быть обработан.

\textbf{Отрицание (вычитание) эффектов.} Система должна давать способ доказать внешнюю
чистоту вычисления: некоторые эффекты используются внутри и не влияют на внешне
наблюдаемое поведение. Например, функция \lstinline|catch| в языке Koka~\cite{
leijen2014koka, leijen2017type} перехватывает исключения переданного вычисления,
оставляя прочие эффекты неизменными:
\[catch : \forall \mu\ap a\ldotp (unit\to^{\{exn|\mu\}} a, exception \to^{\mu} a) \to^\mu a\]
Здесь и далее запись $\{exn|\mu\}$ обозначает множество эффектов, содержащее метку
$exn$ и произвольный <<хвост>> $\mu$; подробно эта нотация обсуждается в
разделе~\ref{sec:rows}.

\textbf{Инкапсуляция эффектов}~\cite{lindley2018encapsulating}. Представим функцию
высшего порядка, которая бросает и ловит исключения для своих внутренних нужд. Без
системы эффектов такая функция может перехватить исключение из вызова переданной ей
функции, тогда как пользователь ожидал его распространения до места вызова. Мы
различаем два аспекта инкапсуляции. Во-первых, система должна предотвращать случайную
обработку, то есть явно отслеживать, что определённые эффекты будут обработаны.
Во-вторых, должен существовать способ замаскировать эффекты-детали реализации от
перехвата ближайшим обработчиком.

\textbf{Подэффектирование (subeffecting).} Эффектную функцию должно быть можно
применять в более разрешающем контексте, чем она требует. Это существенно для гибкости
и удобства, поэтому на типах эффектов нужно отношение подтипизации, что может усложнять
вывод типов.

Эффекты транзитивны вдоль рёбер динамического графа вызовов: эффекты функции включают
эффекты всех вызываемых ею функций~\cite{boruch2023capturing}. Это наблюдение указывает
на главную трудность удобства использования; возможные ответы --- перечислять метки
эффектов в типе, делать набор меток неупорядоченным и полноправным для поддержки
полиморфизма эффектов, обеспечивать надёжный вывод типов эффектов и поддерживать
автоматические пользовательские преобразования эффектов.

Исторически первая система эффектов была предложена для статического обнаружения
ограничений планирования в параллельном программировании выражений с побочными
эффектами~\cite{lucassen1988polymorphic}; тип расширялся метками эффектов и регионами.
Позднее Вадлер и Тиманн показали~\cite{wadler2003marriage}, что программирование с
монадами~\cite{moggi1988computational, wadler1995monads} способно поглотить ранние
системы эффектов. Эту ветвь мы не рассматриваем: для моделирования эффектов в языках
с вызовом по значению монады не обязательны~\cite{liang1995monad, kiselyov2013extensible}. В работе мы сосредоточимся на конструкциях, задуманных скорее практичными, чем максимально
выразительными: нас интересуют системы, реализованные в языках программирования и
рассчитанные на прикладного программиста, а не только на метатеорию. Показательно,
что механизмы определеня собственных эффектов проникают в мейнстрим и без статической дисциплины: OCaml~5 включил обработчики эффектов, вовсе не отслеживая их в
типах~\cite{sivaramakrishnan2021retrofitting}, что делает вопрос о практичной системе
эффектов лишь острее. Прежде чем перейти к конкретным семействам таких конструкций,
сформулируем точку зрения, которая позволит рассматривать их единообразно.


\section{Системы эффектов как управление свободными переменными}
\label{sec:idea}

Опираясь на введённые выше понятия, в этом разделе мы формулируем центральную идею
работы на практических примерах и показываем, что прагматические мотивы систем
эффектов возникают естественно из обычных сценариев программирования.

\subsection{Путь к неявным параметрам}
\label{subsec:implicits}

Рассмотрим сквозной пример: функция бизнес-логики выполняет запрос к хранилищу в
базе данных. В первой версии она сама устанавливает соединение:
\begin{lstlisting}[language=colang]
    fun business1() { let db = Db.open(...); db.query("...") }
\end{lstlisting}
У этого решения есть недостатки. Во-первых, тестируемость ограничена: подменить
соединение макетом непросто. Во-вторых, функция производит наблюдаемые побочные
эффекты --- обращения к базе, --- которые являются частью её внешнего контракта, но
никак не отражены в сигнатуре.

Вторая версия заметно лучше: функция параметризована соединением, и место вызова
может передать конкретную реализацию.
\begin{lstlisting}[language=colang]
    fun business2(db: Db) { db.query("...") }
    fun caller() { let db = Db.open(...); business2(db) }
\end{lstlisting}
Теперь наблюдаемые эффекты явно выражены через параметр \lstinline|db|. Однако за
это приходится платить синтаксическим шумом в месте вызова: административные
усилия по явной передаче параметров растут. Для одного аргумента это мелочь, но если
четыре таких параметра нужно протянуть через несколько уровней стека вызовов,
бремя становится ощутимым. Практика, претендующая на широкое распространение, должна
быть не менее удобной, чем менее качественные альтернативы.

Чтобы сократить рутину, в третьей версии обычный параметр заменяется
\emph{динамически связываемой свободной переменной} \lstinline|db|. Такие переменные
получают своё значение поиском в динамической области видимости вызова.
\begin{lstlisting}[language=colang]
    fun business3() { db.query("...") }
    fun caller() { let db = Db.open(...); business3() }
\end{lstlisting}
Но это достигается ценой утраты других преимуществ: теряются статическая типизация и
явное представление контекстных зависимостей.

Чтобы вернуть статическую безопасность, аппроксимируем динамически связываемые
свободные переменные \emph{неявными параметрами} функции~\cite{lewis2000implicit}. В
месте объявления введём ключевое слово \lstinline|context| для группы неявных
параметров, автоматически подставляемых компилятором на месте вызова; чтобы сделать
значение доступным для разрешения неявных параметров, используем объявление
\lstinline|context let|.
\begin{lstlisting}[language=colang]
    context(db: Db)
    fun business4() { db.query("...") }

    fun caller() {
        context let db = Db.open(...)
        business4()
    }
\end{lstlisting}
В терминах эффектов и контекстов \lstinline|context let| вводит контекст выполнения,
ограниченный соответствующей лексической областью и предоставляющий доступ к
хранилищу. Неявные параметры выражают \emph{требования} функции к её контексту. Чтобы
вызвать такую функцию, контекст должен подтвердить свою пригодность, предоставив
конкретные значения этих параметров. Мы называем такие значения \emph{возможностями}
(capabilities); логически они служат свидетельствами пригодности контекста. Важно,
что до сих пор рассуждение опиралось лишь на механизм связываний и не требовало
глубокого понимания абстрактного понятия эффекта.

Исключения и другие алгебраические операции отслеживаются
аналогично~\cite{odersky2021safer}. Например, конструкцию
\lstinline|throw SomeException()| можно понимать как вызов функции, требующей неявного
параметра типа \lstinline|CanThrow<SomeException>|; конструкции, обрабатывающие
исключения, соответственно предоставляют подходящие значения-свидетельства. Сегодня
ряд распространённых языков поддерживает неявные параметры, и механизмы их вывода
хорошо изучены~\cite{KEEP-context-parameters, oliveira2010type, lewis2000implicit}.

\subsection{Сохранение безопасности эффектов с помощью escape-анализа}
\label{subsec:escape}

Как сформулировано в разделе~\ref{sec:general}, система эффектов должна гарантировать
\emph{безопасность эффектов}: любое вычисление выполняется в контексте, который его
поддерживает. В нашей
постановке контекст предоставляет возможности как свидетельства пригодности, поэтому
эти возможности не должны \emph{утекать} за пределы своих контекстов. Например, если
значение \lstinline|CanThrow| утечёт за пределы блока \lstinline|try-catch|, мы
потеряем гарантию того, что все исключения обработаны, а с ней и безопасность
эффектов:
\begin{lstlisting}[language=colang]
    fun unsafe() {
        var f
        try { // context let h: CanThrow<SomeException>
            let g = () => throw SomeException() // captures h
            f = g // leaks a computation that captured the capability
        } catch (e: SomeException) { ... }
        f() // throws SomeException outside the handler
    }
\end{lstlisting}
Чтобы не допустить утечки возможностей, необходима та или иная форма
escape-анализа~\cite{park1992escape}. Поскольку системы эффектов работают на уровне
типов, естественно применять и escape-анализ на уровне
типов~\cite{xie2022first}; а так как такой анализ должен учитывать поток возможностей
через границы функций, его результаты должны быть отражены прямо в сигнатурах:
сигнатура фиксирует поток данных функции --- какие аргументы используются лишь по
месту, а какие могут быть встроены в результат. Идея делать время
жизни значения объектом типизации восходит к исчислению регионов Тофте и
Талпина~\cite{tofte1997region}; цель та же --- заменить глобальный анализ локальными
условиями на границах функций. Конкретные способы записи такого анализа в типах ---
значения второго класса, явная упаковка, множества захвата --- рассматриваются в
разделе~\ref{sec:capabilities}; здесь нам важно само требование, а не техника.

Важно, что escape-анализ решает задачу \emph{удержания}, не запрещая сам захват:
замыканию можно позволить захватить возможность, если известно, что оно не покинет её
область. Например, лямбде, переданной в энергичную \lstinline|map|, безопасно
обращаться к возможности ввода-вывода из объемлющей области: сигнатура \lstinline|map|
гарантирует, что функциональный аргумент вызывается лишь по месту и не встраивается в
результат, поэтому захваченная возможность не переживёт свой контекст. Здесь
возможность --- лексически видимая свободная переменная лямбды. Способность
замыкаться над возможностями и составляет основу приёма, известного как
\emph{контекстный полиморфизм}~\cite{brachthauser2020effects, brachthauser2022effects}:
функции высшего порядка остаются полностью прозрачными по отношению к эффектам, и не
нужны явные полиморфные переменные, пробегающие строки эффектов. Если же функция,
подобно ленивой \lstinline|map|, встраивает аргумент в результат, её сигнатура должна
честно сообщать об этом --- и тогда ограничения области переносятся на результат.

Польза escape-анализа на уровне типов выходит далеко за пределы построения систем
эффектов: его можно использовать для сокращения числа размещений в куче~\cite{
lorenzen2024oxidizing}, для заимствования при управлении ресурсами на основе
владения~\cite{matsakis2014rust, lorenzen2024oxidizing} и для обеспечения безопасности
нелокальных возвратов~\cite{akhin2021kotlin}.

\subsection{Жизненный цикл контекстного требования}
\label{subsec:idea-summary}

Предыдущее обсуждение по существу велось о свободных переменных. Терм неизбежно
зависит от своего контекста через свободные переменные, приобретая разную семантику в
зависимости от контекста, в котором он вычисляется или определяется. Свободные
переменные --- вот мишень систем эффектов, тогда как связанные регулируются
каноническими системами типов.

Наш взгляд вырастает из простого наблюдения: одно и то же контекстное требование
проходит \emph{жизненный цикл}, а разные системы эффектов делают явной в типах разную
его фазу (рис.~\ref{fig:lifecycle}). Требование \emph{возникает невыполненным}: тело
функции нуждается в чём-то от контекста, но ещё не получило этого. Затем оно
\emph{выполняется} --- контекст предоставляет подходящее значение, которое мы вносим
неявным аргументом. Полученное значение \emph{захватывается} замыканием и должно
\emph{удерживаться} в области своего контекста, что обеспечивает escape-анализ на уровне
типов. Наконец, при пересечении границы контекста захваченное требование может быть
\emph{вновь освобождено} --- переведено обратно в невыполненное и снова отражено в типе.

Двум концам этого цикла отвечают две классические дисциплины разрешения свободных
переменных. Невыполненная фаза динамична по духу: значение требования (обработчик,
возможность) определяется вызывающим контекстом. Выполненная (захваченная) фаза лексична: требование уже
получено и зафиксировано по месту захвата замыканием. Соответствие двух фаз сведено в
табл.~\ref{tbl:free-vars}.

\begin{figure}[t]
\centering
\begin{tikzpicture}[
  node distance=1.0cm and 1.9cm,
  phase/.style={draw, thick, rounded corners, align=center, font=\footnotesize,
                inner sep=4pt, minimum width=2.4cm, minimum height=0.95cm},
  lbl/.style={font=\footnotesize, align=center}]
  \node[phase] (unf) {Невыполнено\\\emph{Unfulfilled}};
  \node[phase, right=of unf] (ful) {Выполнено\\\emph{Fulfilled}};
  \node[phase, right=of ful] (cap) {Захвачено\\\emph{Captured}};
  \draw[thick, -{Latex}] (unf) to[bend left=14] node[lbl, above] {неявный\\параметр} (ful);
  \draw[thick, -{Latex}] (ful) to[bend left=14] node[lbl, above] {захват\\замыканием} (cap);
  \draw[thick, -{Latex}] (cap) to[loop above] node[lbl] {escape-анализ} (cap);
  \draw[thick, -{Latex}] (cap) to[bend left=30] node[lbl, below] {модальность:\\повторное освобождение} (unf);
\end{tikzpicture}
\captionisp{Жизненный цикл контекстного требования. Строковые системы делают явной фазу «невыполнено», системы на основе capabilities --- «выполнено» и «захвачено» с удержанием, модальные --- обратный переход из «захвачено» в «невыполнено»}{Lifecycle of a contextual requirement: row-based systems make the unfulfilled phase explicit, capability-based systems the fulfilled/captured phase, and modal systems the backward transition}
\label{fig:lifecycle}
\end{figure}

\begin{table}[t]
  \captionisp{Две фазы жизненного цикла контекстного требования}{%
  Two phases of a contextual requirement's lifecycle}
  \label{tbl:free-vars}
  \centering
  \begin{tabular}{>{\raggedright\arraybackslash}p{3.3cm}>{\raggedright\arraybackslash}p{3.9cm}>{\raggedright\arraybackslash}p{3.9cm}}
  \toprule
  \textbf{Фаза требования} & \textbf{Невыполнено (динамически)} & \textbf{Захвачено (лексически)} \\
  \midrule
  Смысл                    & ещё не полученное требование к контексту & уже полученное требование, захваченное замыканием \\
  Механизм                 & неявные параметры                 & escape-анализ на уровне типов \\
  Отсюда программист задаёт & дополнительные параметры функции     & поток данных функции \\
  \bottomrule
  \end{tabular}
\end{table}

Одно семейство систем эффектов может задействовать несколько фаз сразу. Так, рассмотренный выше сценарий с
возможностями проходит обе фазы: получение возможности через неявный параметр
(выполнение требования) и её удержание escape-анализом (захват).

Таким образом, систему эффектов можно строить пошагово, как сочетание двух хорошо
спроектированных языковых возможностей --- неявных параметров и escape-анализа на
уровне типов. В
дальнейших разделах мы пользуемся этой оптикой, чтобы заново взглянуть на существующие
семейства систем эффектов.

Оговоримся, что мы предлагаем объединяющую точку зрения, а не формальную
классификацию. Она наиболее естественна для эффектов, возникающих из взаимодействия с
контекстом: алгебраических операций, исключений, доступа к возможностям. Некоторые же
свойства, которые также принято выражать системами эффектов --- расходимость
(незавершаемость), оценки стоимости, метки информационной безопасности, --- удобнее понимать
как предикаты на вычислении; для них прочтение через свободные переменные становится менее
буквальным. В настоящей работе мы сосредоточены на первом классе эффектов, где это прочтение
наиболее точно.


\section{Строковые системы эффектов}
\label{sec:rows}

Большинство систем эффектов представляют набор эффектов \emph{строками}
(rows)~\cite{gaster1996polymorphic}. В этом разделе рассматриваются классические
строковые системы, отличающиеся тем, что хранят в строках сведения обо \emph{всех}
эффектах функций. Метки --- это имена эффектов, возможно применённые к аргументам
типа, а строка есть неупорядоченный набор меток. Например, функция, выполняющая два
эффекта, перечисляет их в типе:
\begin{multline*}
    \lambda xs\ldotp map \ap (\lambda x\ldotp print\ap x\seq yield\ap x) \ap xs \\ : list\ap nat \to^{\{print, ~yield \ap nat\}} list \ap unit
\end{multline*}
Функции высшего порядка производят по меньшей мере те же эффекты, что и их аргументы.
Строковые системы выражают это \emph{полиморфизмом по строкам}~\cite{
gaster1996polymorphic}: переменные типа могут обозначать строки и вклеиваться в другие
строки. Вертикальная черта расширяет строку строковой переменной $\mu$:
\begin{multline*}
    printMap : \forall \mu \ap a \ap b\ldotp (a \to^\mu b) \to list \ap a \to^{\{print|\mu\}} list \ap b
\end{multline*}
В идеале каждая функция высшего порядка должна быть параметризована эффектами своих
функциональных аргументов, что делает плавную миграцию к такой системе эффектов невозможной и
вносит значительное синтаксическое загрязнение.

Пример строковой системы --- проверяемые исключения Java~\cite{gosling2000java}: язык
статически удостоверяется, что все проверяемые исключения пойманы или явно объявлены в
сигнатуре метода. Эта черта печально известна нелюбовью
программистов~\cite{checked-exceptions}: строка исключений --- специальная (ad hoc)
конструкция, которой нельзя параметризовать функцию, поэтому написать метод высшего
порядка, пробрасывающий произвольное число проверяемых исключений, невозможно.

Многие исследовательские языки реализуют классические строковые системы. При этом
необходимость подтипизации и механизма вычитания обработанных эффектов обычно ведёт к
переусложнению при попытке уложиться в стиль вывода типов
Хиндли--Милнера~\cite{hindley1969principal, milner1978theory}. Так, система языка
Links поддерживает аннотации и полиморфизм присутствия меток в
строке~\cite{hillerstrom2016liberating}. Мы же сосредоточимся на системе языка Koka,
простой, но достаточно мощной~\cite{leijen2014koka, leijen2017type}.

Koka различает дублирующиеся метки в строках, $\{l,l\}\neq \{l\}$, следуя идее
меток с областью действия (scoped labels) для
записей~\cite{leijen2005extensible}. Это позволяет
сделать формы обработки прямолинейными и при этом дать обработчику пользоваться
исключениями. Koka предотвращает случайную обработку, выводя подходящий тип, а операция
\lstinline|mask| скрывает эффект от ближайшего обработчика, инкапсулируя его:
\begin{lstlisting}[language=koka]
    fun catch(f : () -> <exn|e> a,
              h : err -> <|e> a) : <|e> a

    fun test(f : () -> <|e> int) : <|e> int
      catch(fn () mask<exn> { f() }, handle)
\end{lstlisting}
Типизация эффектов Koka укладывается в стиль Хиндли--Милнера, а подтипизация строк
обеспечивается добавлением дополнительных полиморфных переменных $\{exn|\mu\}$.

Когда строки не поддержаны напрямую, их эмулируют другими средствами. Библиотека
\texttt{mtl} в Haskell\footnote{\url{https://hackage.haskell.org/package/mtl}}
перечисляет эффекты ограничениями классов типов~\cite{jones1995functional}: будучи
неупорядоченными и поддерживая подтипизацию, они выступают как ad hoc строки.
Библиотеки обработчиков эффектов в Haskell тоже используют классы типов для открытых
объединений~\cite{swierstra2008data}; ещё один подход --- эмулировать строки
типами-пересечениями~\cite{xie2020row}, практическим примером чего служит библиотека
ZIO\footnote{\url{https://zio.dev/}} для Scala.

Классические строковые системы --- самый прямолинейный подход, поэтому новые языки
часто берут его за основу: так, в языке Flix вывод полиморфных эффектов опирается на
булеву унификацию~\cite{madsen2020polymorphic}, а язык Unison реализует строковую
дисциплину в форме \emph{способностей} (abilities)~\cite{unison-abilities}. Однако такие
системы отслеживают \emph{слишком много}
сведений, что ведёт к избыточной многословности.

\subsection{Прочтение: типизация динамически видимых свободных переменных}
\label{subsec:rows-reading}

С точки зрения нашей перспективы строковая система делает явной \emph{невыполненную}
фазу цикла --- сами операции-эффекты как свободные переменные с динамической
дисциплиной, --- в эффектном стиле, аналогичном монаде Reader. Строка эффектов на стрелке
есть в точности перечень \emph{невыполненных контекстных требований} тела функции:
каждая метка \texttt{print}, \texttt{yield} или \texttt{exn} --- это свободная
переменная-операция, значение которой (обработчик) будет найдено в динамической области
вызова.

Этот взгляд объясняет и сильные, и слабые стороны семейства. Поскольку отслеживаются
именно невыполненные требования, обработка эффекта (предоставление обработчика)
естественно \emph{вычитает} метку из строки --- это и есть отрицание эффектов из
раздела~\ref{sec:general}. С другой стороны, так как требования динамические,
каждое из них приходится протаскивать через тип каждой промежуточной функции; отсюда
многословность и необходимость повсеместного полиморфизма по строкам.

Как отмечено в разделе~\ref{subsec:idea-summary}, «динамичность» здесь --- о духе
дисциплины (связывание разрешается вызывающим контекстом), тогда как сама проверка
требований остаётся \emph{статической}; именно это отличает строковые системы от
нетипизированного динамического связывания.

% Отождествление строк с динамической дисциплиной имеет и известные границы, которые
% скорее подтверждают предлагаемое прочтение, чем опровергают его. Во-первых, существуют
% строковые системы с \emph{лексически} связываемыми обработчиками: связь операции с
% обработчиком фиксируется по месту определения, а не поиском в динамическом контексте
% вызова~\cite{biernacki2019binders}; туннелирование обработчиков аналогично сочетает строковую
% запись с лексической дисциплиной, решая проблему случайной
% обработки~\cite{zhang2016accepting, zhang2020handling}. Во-вторых, техника передачи
% свидетельств (evidence passing) компилирует строковые эффекты Koka в явную передачу
% значений-свидетельств обработчиков~\cite{xie2021generalized} --- то есть невыполненные
% требования систематически преобразуются в выполненные, полученные как параметры. Оба
% наблюдения --- прямая опора нашего тезиса: фазы табл.~\ref{tbl:free-vars} суть два
% представления одной и той же информации, связанные механическим переводом, а выбор
% семейства --- вопрос того, какую фазу язык делает исходной.


\section{Системы на основе capabilities}
\label{sec:capabilities}

Ранее мы видели эффекты как вызовы <<особых>> функций, например \texttt{print}, и
перечисляли соответствующие метки в типах, что ведёт к многословным системам:
\[\lambda x\ldotp print\ap x\seq x + 1 : nat \to^{\{print\}} nat\]
Теперь взглянем на отслеживание эффектов иначе. Предположим, что в программе есть
<<особые>> объекты --- \emph{возможности} (capabilities); эффекты выполняются только
через них, поэтому система может отслеживать эти объекты вместо самих операций.
Возможности можно понимать как выданные контекстом разрешения на выполнение
эффектов:
\[\lambda x\ldotp console.print\ap x\seq x + 1 : nat \to^{\{console\}} nat\]
Функция получает возможности несколькими путями. Во-первых, она может принимать их как
аргументы; тогда о возможностях рассуждают так же, как о связываниях~\cite{
brachthauser2022effects}, --- по существу заменяя тип-эффект явным параметром:
\[\lambda \ap console \ap x \ldotp console.print\ap x\seq x + 1 : Console \to nat \to nat\]
Явная передача возможностей аргументами была бы громоздкой (та самая проблема, которую
мы решаем эффектами), поэтому язык должен поддерживать неявную передачу параметров.
Именно здесь, как отмечает~\cite{boruch2023capturing}, систему эффектов уместно видеть как
средство отслеживания свободных, динамически связываемых переменных; это согласуется с
картиной из раздела~\ref{sec:idea}: система типов контролирует связанные переменные, а
система эффектов --- свободные.

Во-вторых, возможности могут захватываться функциями первого
класса. Так возможности дают лёгкую альтернативу классическому полиморфизму эффектов ---
\emph{контекстный полиморфизм}~\cite{brachthauser2022effects}: функции \texttt{map}
можно дать привычный тип без какого-либо явного полиморфизма по эффектам:
\[map : \forall a \ap b\ldotp (a \to b) \to list \ap a \to list \ap b \]
Наконец, возможны и другие источники возможностей --- это проектные решения; например,
глобальные определения могут использоваться как возможности для упрощения
миграции~\cite{brachthauser2022effects, boruch2023capturing}, а естественным источником
служат и обработчики эффектов~\cite{brachthauser2020effects}.

Системы на основе capabilities можно эмулировать и без встроенной поддержки, в
определённом стиле программирования: функции не пользуются глобальными определениями
для эффектов, а принимают объекты-возможности (службы) и вызывают их методы. Идею можно
показать на Scala с неявными параметрами~\cite{odersky2004scala}; язык развивается в
сторону встроенной проверки исключений с <<магической>> возможностью
\lstinline[language=scala]|CanThrow|~\cite{odersky2021safer}.

К сожалению, обеспечение безопасности систем на основе capabilities требует
дополнительных усилий. Система безопасна, когда она не даёт возможностям утекать за
пределы выделенной области. Самый наглядный пример --- исключения: их нельзя бросать вне
соответствующего \lstinline[language=scala]|try-catch|, однако возможность-исключение
легко утекает. Аналогичные примеры есть и для других эффектов: дескриптор файла,
по существу являющийся возможностью работы с файлом, со временем становится
недействительным. Видна близость понятий ресурса и возможности: у обоих ограниченное
время жизни.

Поэтому нужно различать значения первого класса, передаваемые свободно, и значения
второго класса с ограниченной областью. Одно из решений --- классифицировать значения
на уровне типов и ограничить значения второго класса
правилами~\cite{osvald2016gentrification}: возможности суть значения второго класса;
функции первого класса не могут ссылаться на значения второго класса через свободные
переменные; функции могут возвращать только значения первого класса; в полях объектов и
изменяемых переменных можно хранить лишь значения первого класса. Этот подход довольно
негибок: каррированная \texttt{map} попыталась бы вернуть функцию второго класса, что
запрещено; трудности возникают и с наследованием~\cite{osvald2016gentrification}.

Возможности не обязаны быть явными значениями. Язык
Effekt\footnote{\url{https://effekt-lang.org/}} использует неявную передачу
возможностей на уровне исходного кода, транслируя её в исчисление с явной передачей и
лексически видимыми обработчиками~\cite{brachthauser2020effects}. Effekt тоже использует
строки, но с иной семантикой, чем в разделе~\ref{sec:rows}: строки уже не обозначают все
эффекты контекста, а читаются как требования и предоставления возможностей в сигнатуре:
\begin{lstlisting}[language=effekt]
    def buildString(indent: Int)
      { f: () => () / {Yield[String]} } // block argument
      : String / {Format}               // buildString provides / requires
\end{lstlisting}
Здесь инкапсуляция эффектов работает по умолчанию: трансляция, управляемая типами, не
передаёт возможности функциональному аргументу, если он явно их не запрашивает.

С неявными возможностями мы беспокоимся не об утечке самих объектов-возможностей, а об
отложенных вычислениях, захватывающих возможности. Effekt решает это, делая все функции
блоками второго класса~\cite{brachthauser2020effects}. Но поддержка только блоков
второго класса слишком ограничительна, поэтому был предложен явный механизм
\emph{упаковки}: \texttt{box} делает неявные требования терма явными, а \texttt{unbox}
заставляет систему проверить, что упакованное значение применимо в текущем
контексте~\cite{hannan1998type, brachthauser2022effects}. В правилах ниже суждение
$\Gamma \vdash b : \sigma ~|~ C$ читается так: терм $b$ имеет тип $\sigma$ и захватывает
множество возможностей $C$ (его \emph{множество захвата}); а тип
$\sigma~\mathbf{at}\ap C$ --- это тип упакованного значения, замкнувшегося над $C$. Тогда
\texttt{box} переносит захват из суждения в тип, оставляя сам терм чистым ($\{\}$), а
\texttt{unbox} проверяет, что все возможности из $C$ доступны в текущем контексте:
\begin{gather*}
    \infer[BoxIntro]{\Gamma \vdash \mathbf{box}\ap b : \sigma~ \mathbf{at}\ap C ~|~ \{\}}{\Gamma \vdash b : \sigma ~|~ C}
    \qquad
    \infer[BoxElim]{\Gamma \vdash \mathbf{unbox}\ap e : \sigma ~|~ C}{\Gamma \vdash e : \sigma ~\mathbf{at}~C ~|~ \{\}}
\end{gather*}
Поскольку типы упакованных значений снова перечисляют все эффекты, нужна форма
полиморфизма эффектов; Effekt вводит \emph{полиморфизм по возможностям}, выражающий, что
упакованное значение замыкается над теми же возможностями, что и заданное
значение~\cite{brachthauser2022effects}.

В недавних работах по Scala исследуется \emph{отслеживание захвата} (capture checking),
которое для каждого значения отслеживает, какие возможности оно
захватывает~\cite{boruch2023capturing}. Подход лучше согласуется с
подтипизацией и не разделяет значения на классы столь строго. Каждому значению типа
\lstinline[language=scala]|T| можно приписать тип вида \lstinline[language=scala]|T^{cap}|,
указывающий, что значение --- возможность с ограниченным использованием (его нельзя
вернуть из функции, им нельзя инстанцировать обобщённые типы). Поскольку поддержан
контекстный полиморфизм, функциональные типы тоже являются возможностями. Множества
захвата образуют иерархию подзахвата ($C_1 <: C_2$, если $C_1$ покрывается $C_2$):
например, если \lstinline[language=scala]|f| захватывает только \texttt{c1} и
\texttt{c2}, то $\{f\} <: \{c1, c2\}$. Синтаксическое загрязнение при этом разумно:
бремя возникает лишь при выпуске возможностей или вычислений за пределы области, в
остальном работает вывод. Тем не менее типы, ссылающиеся на переменные-термы (лёгкая
форма зависимых типов), непривычны для среднего программиста, что может затруднить
принятие.

\subsection{Прочтение: лексически видимые свободные переменные и их удержание}
\label{subsec:caps-reading}

С точки зрения нашей перспективы системы на основе capabilities делают явной
\emph{захваченную} фазу цикла. Это семейство проходит обе соседние фазы: возможность
сначала \emph{получается} из контекста (как аргумент, как глобальное определение или из
обработчика) --- по сути через неявный параметр, --- а затем \emph{захватывается}
замыканием как лексически видимая свободная переменная. В типы выносится именно
захваченное, уже выполненное требование из табл.~\ref{tbl:free-vars}; получение же
требования (динамическая фаза) обычно доверяется выводу. Поэтому прочтение не
противоречит тому, что возможности отслеживаются как выполненные, хотя и добываются
динамически.

Отсюда два следствия, объясняющие устройство семейства. Во-первых, контекстный
полиморфизм --- прямое следствие захвата лексически видимых возможностей: функция высшего
порядка не обязана быть полиморфной по эффектам, потому что её аргумент уже несёт в себе
свои возможности. Во-вторых, главной трудностью становится не учёт требований, а
\emph{escape-анализ}: нужно гарантировать, что захваченная возможность не покинет
область своего обработчика. Все рассмотренные механизмы --- значения второго класса,
упаковка \texttt{box}/\texttt{unbox} и отслеживание захвата --- суть разные ответы на
один и тот же вопрос escape-анализа, отличающиеся ценой: жёсткое разделение на классы,
явные аннотации упаковки или множества захвата в типах. Это в точности требование
escape-анализа из раздела~\ref{subsec:escape}, воплощённое тремя разными способами;
корневая возможность \lstinline[language=scala]|cap| играет при этом роль универсального
признака отслеживаемости~\cite{boruch2023capturing}.


\section{Модальные типы эффектов}
\label{sec:modal}

Capabilities позволили описывать лишь часть окружающего контекста --- только те
возможности, что непосредственно требуются вычислению. Главной трудностью был учёт
неявно захваченных возможностей, где все подходы оказывались слишком ограничительными
или требовали зависимых типов. В этом разделе рассматривается система, которая тоже
стремится явно указывать как можно меньше частей контекста, но остаётся целиком на
уровне типов.

Изначально язык Frank предложил \emph{способности} (abilities) и
\emph{подстройки} (adjustments) --- конструкции второго класса как синтаксический сахар
для полиморфизма по строкам~\cite{lindley2017do, convent2020doo}. Позднее была
предложена самостоятельная модальная система, делающая эти конструкции
полноправными~\cite{tang2025modal} и следующая идее использовать контекстную модальную
теорию типов для отслеживания алгебраических эффектов~\cite{zyuzin2021contextual}.
Ключевая мысль: конструкторы-модальности должны помечать \emph{переходы} между
контекстами эффектов, а не повторять весь контекст.

Система использует два конструктора-модальности, параметризуемые метками эффектов; их
значения полноправны (могут храниться и передаваться свободно). Общая идея ---
упаковывать вычисления и приписывать им модальность, описывающую требуемый контекст; в
подходящем контексте вычисления автоматически распаковываются. \emph{Абсолютная}
модальность (квадратные скобки) описывает весь требуемый контекст: например,
\texttt{map} не накладывает ограничений на контекст, а \texttt{gen} требует эффект
\texttt{yield}:
\begin{lstlisting}[language=koka]
    map : []((a -> b) -> List a -> List b)
    gen : [yield](List Int -> 1)
\end{lstlisting}
Пустые скобки \texttt{[]} в типе \texttt{map} говорят, что тело замкнуто и не зависит от
контекста эффектов, поэтому его можно переносить куда угодно; \texttt{[yield]} у
\texttt{gen} --- что для запуска нужен контекст, обрабатывающий \texttt{yield}.
\emph{Относительная} модальность (угловые скобки) описывает относительные требования:
упакованное вычисление может использовать как эффекты окружающего контекста, так и
перечисленные в типе. Так, \texttt{asList} обрабатывает \texttt{yield} и принимает
вычисление, которое среди прочего может выполнять \texttt{yield}:
\begin{lstlisting}[language=koka]
    asList : <yield>(1 -> 1) -> List Int
\end{lstlisting}
Угловые скобки \texttt{<yield>} означают, что аргумент вправе выполнять \texttt{yield}
\emph{в дополнение} к эффектам окружающего контекста, а сам \texttt{asList} этот
\texttt{yield} обрабатывает. Две конструкции преобразуют контекст: обработчики добавляют
эффекты, маскирование --- вычитает. Относительные модальности имеют общий вид
\texttt{<L|R>}, где левая часть \texttt{L} --- строка маскируемых эффектов, а правая
\texttt{R} расширяет окружающий контекст; так, \texttt{<yield>} выше ничего не маскирует
(левая часть пуста) и лишь добавляет \texttt{yield}. Есть правило: переменную
следует использовать в контексте, совместимом с тем, в котором она была связана; оно
обеспечивает и безопасность, и инкапсуляцию эффектов. Полученная система достаточно
выразительна, чтобы кодировать программы строковых систем, пока переменные эффектов на
стрелках ссылаются на ближайшую лексически. Однако она по-прежнему требует множества
явных аннотаций --- упаковки, маскирования и <<замораживания>>~\cite{emrich2020freezeml},
--- а система видов (kinds) дополнительно усложняет картину.

\subsection{Прочтение: повторное освобождение лексически видимых переменных}
\label{subsec:modal-reading}

С точки зрения нашей перспективы модальные типы занимают особое положение: они делают
явными не отдельную фазу, а \emph{переходы} между фазами --- рёбра
рис.~\ref{fig:lifecycle}. Разведём два понятия, смешение которых и создаёт видимость
противоречия. Модальность, приписанная упакованному вычислению, выражает его
\emph{остаточное невыполненное} требование к контексту: в подходящем контексте вычисление
автоматически распаковывается и становится готовым к исполнению. Система же видов
отслеживает нечто иное --- есть ли у значения уже \emph{захваченные} (лексические)
зависимости от контекста. На стыке этих двух и работают модальные системы: значение с
захваченной зависимостью, покидая свой контекст, снабжается модальностью, которая вновь
превращает эту зависимость в невыполненное требование в типе.

Ключевую роль играет система видов: она различает вычисления, которые могут захватывать
лексически видимые свободные переменные (вид $Any$), и те, что заведомо не зависят от
окружающего контекста эффектов (вид $Abs$)~\cite{tang2025modal}. Значения абсолютных
модальностей имеют вид $Abs$ и переносятся между контекстами свободно. Значение же вида
$Any$ может нести в себе захваченную зависимость от контекста: например, функция,
определённая в области действия обработчика \texttt{yield} и использующая этот эффект,
захватывает его как лексически видимую свободную переменную. Когда такое значение
переносится в другой контекст (например, возвращается из обработчика), система не
отвергает программу, а дополняет его тип относительной модальностью, вновь делающей
зависимость явной. Этот механизм мы называем \emph{повторным освобождением}
(uncapturing) переменной: выполненное контекстное требование, захваченное замыканием,
переводится обратно в невыполненное, которое далее отслеживается в типе.

Таким образом, здесь escape-анализ не \emph{запрещает} утечку, как в
разделе~\ref{sec:capabilities}, а \emph{контролирует} её. Подчеркнём, что предотвращение
небезопасных утечек никуда не исчезает: его обеспечивают правило совместимости контекстов
использования и связывания переменной вместе с системой видов; отличие от capabilities
в том, что этот механизм не образует отдельной подсистемы с собственными аннотациями, а
встроен в дисциплину модальностей. Это объясняет и место модальных систем: они работают с
той же захваченной фазой, что и capabilities (табл.~\ref{tbl:free-vars}), но выносят в
типы сам \emph{переход} между фазами, а не запрет утечки, --- отсюда контроль вместо
запрета.


\section{Сравнительный анализ}
\label{sec:comparison}

Сведём три семейства в единую картину. Предложенная точка зрения даёт единую ось
сравнения --- \emph{фазу жизненного цикла контекстного требования}, которую система делает
явной в типах (невыполненную, захваченную или переход между ними). Всё остальное --- что
именно отслеживается, какой требуется полиморфизм и, в особенности, где сосредотачиваются
аннотации --- оказывается \emph{следствием} этого выбора, а не независимым измерением.
Сравнение приведено в табл.~\ref{tbl:compare}.

\begin{table}[t]
  \captionisp{Три семейства систем эффектов вдоль единой оси фаз}{%
  Three families of effect systems along the single phase axis}
  \label{tbl:compare}
  \centering
  \footnotesize
  \begin{tabular}{>{\raggedright\arraybackslash}p{2.0cm}>{\raggedright\arraybackslash}p{2.7cm}>{\raggedright\arraybackslash}p{2.7cm}>{\raggedright\arraybackslash}p{2.7cm}}
  \toprule
  & \textbf{Строковые} & \textbf{Capabilities} & \textbf{Модальные} \\
  \midrule
  Явная фаза цикла & невыполнено & захвачено & переход между фазами \\
  Проходит фазы & невыполнено & выполнено, захвачено & захвачено, невыполнено \\
  Что отслеживается & метки операций в строке & объекты-возможности & переходы между контекстами \\
  Полиморфизм & по строкам & контекстный & по видам/модальностям \\
  Аннотации (следствие фазы) & в сигнатурах всех функций высшего порядка & на границах утечки возможностей & в точках переноса между контекстами \\
  Безопасность эффектов & вычитанием меток & escape-анализом (запрет утечки) & переносом модальности (контроль утечки) \\
  Цена & многословность & сложность escape-анализа & аннотации, система видов \\
  Примеры систем & Koka, Links, проверяемые исключения Java & Effekt, отслеживание захвата в Scala & Frank; модальные типы эффектов~\cite{tang2025modal} \\
  \bottomrule
  \end{tabular}
\end{table}

Картина обнаруживает закономерность. Строковые системы делают явной невыполненную фазу:
они учитывают \emph{невыполненные} требования и потому платят многословностью и
повсеместным полиморфизмом по строкам, зато обработка эффекта прямолинейно вычитает метку.
Системы на основе capabilities и модальные делают явной захваченную фазу: требование уже
выполнено и захвачено замыканием, благодаря чему появляется контекстный полиморфизм и
исчезает нужда в явных переменных эффектов, но центральной проблемой становится
escape-анализ. Capabilities решают его \emph{запретом} утечки (значения второго класса,
упаковка, множества захвата), а модальные системы --- \emph{контролем} утечки через
повторное освобождение переменных.

Конкретность различия удобно увидеть на одной и той же функции \texttt{map}
(рис.~\ref{fig:map-three}). В строковой системе её тип несёт полиморфную переменную
эффекта (в Koka: \lstinline[language=koka]|fun map(xs: list<a>, f: (a) -> e b): e list<b>|);
в системе на основе capabilities эффекты в типе не упоминаются вовсе благодаря контекстному
полиморфизму; в модальной системе тело оборачивается абсолютной модальностью. Один и тот же
программный объект получает три разные сигнатуры, и различие в точности отражает, какую
фазу цикла делает явной система.

\begin{figure}[t]
\centering
\begin{tabular}{@{}ll@{}}
\toprule
\textbf{Семейство} & \textbf{Тип \texttt{map}} \\
\midrule
Строковая     & $(a \to^{\mu} b) \to list\ap a \to^{\mu} list\ap b$ \\
Capabilities  & $(a \to b) \to list\ap a \to list\ap b$ \\
Модальная     & $[\,]((a \to b) \to List\ap a \to List\ap b)$ \\
\bottomrule
\end{tabular}
\captionisp{Одна функция \texttt{map} в трёх семействах: строковая несёт полиморфную переменную эффекта $\mu$, capabilities не упоминают эффекты благодаря контекстному полиморфизму, модальная оборачивает тело абсолютной модальностью}{The single function map across the three families}
\label{fig:map-three}
\end{figure}

Другой поучительный пример --- инкапсуляция эффекта-детали реализации из
раздела~\ref{sec:general}: функция внутренне бросает и ловит исключение, оставаясь
внешне чистой. В строковой системе обработчик вычитает метку \texttt{exn} из строки, а
от случайного перехвата чужих исключений защищает маскирование. В системе на основе
capabilities достаточно того, что возможность \lstinline|CanThrow| вводится и
используется локально: escape-анализ удостоверяет, что она не покидает область
обработчика, и снаружи эффект не виден. В модальной системе обработчик снимает
соответствующую компоненту модальности, а правило совместимости контекстов не даёт
перехватить исключения, пришедшие из другого контекста. Один и тот же контракт внешней
чистоты достигается тремя дисциплинами учёта одной и той же свободной переменной.

Предложенное прочтение связано с теорией коэффектов~\cite{petricek2014coeffects}:
невыполненная фаза соответствует эффектному учёту (со стороны производимых вычислением
запросов), а захваченная --- коэффектному (со стороны получаемых из контекста ресурсов).

\subsection{Развязка: система эффектов из двух знакомых частей}
\label{subsec:synthesis}

Единая ось подсказывает конструктивную программу: систему эффектов можно собрать из двух
независимо полезных и уже знакомых практикам возможностей --- \emph{неявных параметров}
(для невыполненной фазы) и \emph{escape-анализа на уровне типов} (для захваченной). Обе
половины уже проиллюстрированы выше на сквозных примерах, и их достаточно.

Первая половина отвечает за выполнение требования. Функция \lstinline|business4| из
раздела~\ref{subsec:implicits} не открывает соединение сама и не тянет его лишним
аргументом: она объявляет \lstinline|context(db: Db)|, а место вызова предоставляет
возможность через \lstinline|context let|. Так невыполненное контекстное требование
разрешается неявным параметром --- ровно то, что делают строковые системы, только явная
метка эффекта заменена значением-свидетельством.

Вторая половина отвечает за безопасность. Как только возможность получена и может быть
захвачена замыканием (раздел~\ref{subsec:escape}), нужно гарантировать, что она не
переживёт свой контекст: лямбда, отданная энергичной \lstinline|map|, вправе замкнуться
над возможностью, потому что сигнатура \lstinline|map| обещает вызвать её лишь по месту;
а ленивая \lstinline|map|, встраивающая аргумент в результат, обязана честно перенести
ограничение области на результат. Это и есть escape-анализ на уровне типов, который вдобавок
бесплатно даёт контекстный полиморфизм.

Собранные вместе, эти две части и образуют систему эффектов: неявные параметры вводят и
разрешают требования, а escape-анализ удерживает уже разрешённые. Три рассмотренных
семейства оказываются разными способами сделать явной ту или иную фазу этого механизма, а
не принципиально различными конструкциями. Практический выигрыш в том, что многие массовые
языки уже поддерживают хотя бы одну из двух половин, поэтому отслеживание эффектов можно
внедрять постепенно.


\section{Заключение}
\label{sec:conclusion}

Мы показали, что статический контроль над эффектами важен для надёжности программ и
построения ясных границ абстракций. Главный тезис работы состоит в том, что три, на
первый взгляд, несвязанных семейства систем эффектов --- строковые, на основе
capabilities и модальные --- суть варианты одного механизма: управления свободными
переменными вычисления на уровне типов. Контекстное требование проходит единый жизненный
цикл --- возникает невыполненным, выполняется значением из контекста через неявный
параметр, захватывается замыканием и удерживается escape-анализом, а на границе контекста
может быть освобождено вновь, --- и семейства различаются лишь тем, какую фазу этого цикла
они выносят в типы. Строковые системы делают явной невыполненную фазу; системы на основе
capabilities отслеживают уже захваченные возможности и не дают им утекать; модальные
системы помечают переходы между фазами, возвращая захваченные требования в невыполненное
состояние и тем самым контролируя утечку.

Этот единый взгляд не только упорядочивает существующие конструкции, но и подсказывает
конструктивный путь к более простым и практичным системам эффектов --- собрать их из
неявных параметров и escape-анализа на уровне типов. Мы надеемся, что системы эффектов
достигнут необходимой простоты и сократят число вспомогательных церемоний, требуемых для
их принятия в массовых языках программирования.

\subsection*{Ограничения и направления дальнейшей работы}

У предложенной перспективы есть ряд ограничений. Во-первых, проведённое сравнение качественное;
количественная оценка нагрузки от аннотаций --- например, сопоставление объёма аннотаций для
одной и той же программы в трёх семействах --- выходит за рамки работы и составляет предмет
дальнейшего исследования. Во-вторых, охват сознательно сужен до трёх семейств: градуированные
(graded) и индексированные системы, а также эффекты высшего порядка (scoped
effects)~\cite{wu2014effect} затронуты лишь вскользь, хотя и они допускают прочтение в
терминах свободных переменных; в частности, градуированные модальные типы языка
Granule~\cite{orchard2019quantitative} родственны конструкциям раздела~\ref{sec:modal},
но отслеживают количественные сведения об использовании контекста, а не переходы между
контекстами. В-третьих, как отмечено в разделе~\ref{sec:idea},
предлагаемое прочтение менее точно для эффектов, не связанных с взаимодействием с контекстом
(расходимость, стоимость вычисления). Наконец, главный открытый инженерный вопрос ---
объединение двух дисциплин (неявных параметров и escape-анализа) в одном практичном языке с
приемлемым выводом типов и обратной совместимостью; это направление дальнейшей работы, а
развитые модели обработчиков эффектов~\cite{van2022handling, sivaramakrishnan2021retrofitting}
показывают, что необходимая для этого среда исполнения постепенно становится частью
мейнстрима.


\maketitleen

\begin{otherlanguage}{english}
\printbibliography
\end{otherlanguage}

{\vskip 12pt\normalfont\sffamily\bfseries\large Информация об авторах / Information about authors}
\setlength{\parskip}{6pt}

Андрей Сергеевич СТОЯН --- аспирант департамента информатики Национального
исследовательского университета <<Высшая школа экономики>> (Санкт-Петербург). Сфера
научных интересов: проектирование языков программирования, системы типов, системы
эффектов.

Andrey Sergeevich STOYAN --- PhD student at the Department of Informatics, National
Research University Higher School of Economics (Saint Petersburg). Research interests:
programming language design, type systems, effect systems.

\end{document}
